\section{Compared with Other Methods}
In the context of this discussion, we assume that the NBA for an LTL specification $\phi$ has already been constructed. The translation from an LTL formula to an NBA involves an exponential complexity with respect to the size of the formula \cite{vardi2005automata}. The complexity of constructing the complement of an NBA is known to be EXPTIME \cite{complementnba}.
% On the other hand, constructing the complement of an LTL formula is linear time.
% 下面对比的三种方法中使用的都是state-based NBA,
\subsection{The state-triplet Approach}
Barrier certificate (\ref{SCDEF}) separates the safe region $\mathcal{X}_{safe}$ from the unsafe region $\mathcal{X}_{u}$, ensuring that the system state remains within $\mathcal{X}_{safe}$ at all times and thus preventing the system behavior from visiting $\mathcal{X}_{u}$. If consider $\mathcal{X}_{VF}$ as the unsafe set, finding the corresponding barrier certificate $B$ can guarantee the persistence of the system with respect to $\mathcal{X}_{VF}$. The state-triplet approach leverages barrier certificates to cut off all simple paths from the initial state \(q_0\) to any accepting state in the NBA automaton representing the complement of the LTL specification $\phi$. If one can find all certificates for all simple paths, this method ensures that runs of \emph{words} in $\mathfrak{L}(\mathfrak{S})$ never visit any accepting state, consequently, \(\mathfrak{S} \models_{\mathfrak{L}} \phi\).

The concept of reachability is consistent between the \emph{state safety certificate} and the state-triplet approach, but the \emph{state safety certificate} achieves this by leveraging barrier certificates on the product system $\mathfrak{S} \times \mathcal{A}$. While the state-triplet needs to block all simple paths leading to accepting states from $q_0$, the number of certificates in the state-triplet method is related to the number of triplets. In addition, by finding a state safety certificate, it can avoid the need to unroll the automaton and expand all simple paths and then block them individually. The \emph{transition persistence certificate} proposed in this paper allows runs of \emph{words} in $\mathfrak{L}(\mathfrak{S})$ to visit accepting states but ensures that they visit accepting transitions finitely often by enforcing a decrease in the certificate value.

% So combination of state safety certificates and transition persistence certificate, 在保证原有能力的同时，也能验证简单路径可达的并且有限访问的性质验证，减缓了三元组方法的保守性。

The combination of state safety certificates and transition persistence certificates can reduce the conservatism of the state-triplet method. While maintaining its original capability of proving that a system satisfies specifications by cutting off reachability, this combination can also verify systems where there exist runs that can reach accepting states from $q_0$ and meet the specifications.

\subsection{Closure Certificate}
According to the definition of the LTL closure certificate (\ref{LTLCC}), the certificate is defined on $\mathcal{X} \times Q \times \mathcal{X} \times Q$. The increase in state dimensionality leads to longer synthesis times in practice. Therefore, searching for an LTL closure certificate imposes a heavy computational burden. For our certificates, whether it is the transition persistence certificate defined on $\mathcal{X}$ or the state safety certificate defined on $\mathcal{X} \times Q$, the state space is smaller compared to that of the closure certificate. Additionally, the number of transition persistence certificates is related to $|Acc^{pair}|$, and the number of state safety certificates is related to $|Acc^{start}|$. Thus, the search space for all certificates is at most $\mathcal{X} \times Q \times Q$. Subsequent case studies will demonstrate that for cases previously involving closure certificates \cite{murali2024closure}, the state safety certificates and transition persistence certificates can be synthesized in seconds, resulting in significant improvement in synthesis efficiency.
